\section{Related Work}
\label{sec:related}

We organise the literature around a single question: when a population method
attributes its performance to structure, what evidence is offered, and what
would settle the matter? Section~\ref{sec:rel:topology} covers methods in
which topology is a design choice, \ref{sec:rel:measure} those in which it is
a measured object, \ref{sec:rel:theory} the theory that describes what such
updates do, \ref{sec:rel:adaptive} adaptive components and their credit
assignment, and \ref{sec:rel:method} the methodological literature on
attribution.

\subsection{Topology as a design choice}
\label{sec:rel:topology}

That the interaction graph shapes a population's behaviour is among the
better-established empirical facts in the field. Kennedy and Mendes varied the
particle swarm's communication topology across a large set of graphs and found
performance to depend on it systematically, with denser neighbourhoods
favouring rapid convergence and sparser ones preserving
diversity~\cite{kennedy2002population}; the fully informed variant~\cite{mendes2004fully}
and comprehensive learning~\cite{liang2006clpso} pursue the same insight through
the update rule, and decentralised differential evolution through the
population graph~\cite{dorronsoro2011cellular}. Cellular evolutionary algorithms make
the mechanism explicit: Alba and Dorronsoro parameterise the relationship
between neighbourhood and grid by a single ratio, show that this ratio governs
selection pressure, and schedule it during the run to move between exploration
and exploitation~\cite{alba2005exploration}. Takeover-time analysis supplies
the accompanying theory, predicting the rate at which the population collapses
onto an incumbent as a function of neighbourhood shape.

This line establishes (P1) in the terminology of
Section~\ref{sec:claim}---heterogeneity and connectivity matter---but it makes
no claim of the form (P2). The grid is fixed a~priori and no node is asserted
to be more informative than another by virtue of its position.

A more recent family does make such a claim. Here a structural statistic
computed on the population graph weights the update directly, so that agents
learn preferentially from neighbours that occupy distinguished positions. The
graphs are typically Watts--Strogatz small worlds~\cite{watts1998collective}
and the statistic typically betweenness, computable in $O(N\!\cdot\!M)$ by
Brandes' algorithm~\cite{brandes2001faster}. The
optimizer we adopt as our primary case study places candidate solutions on a
small-world graph and weights influence by betweenness centrality, arguing
that high-betweenness agents bridge regions of the search
space~\cite{jalali2026social}. Related proposals substitute degree,
closeness, eigenvector centrality or PageRank. Our concern throughout is not
whether these methods perform competitively---several do---but whether the
structural statistic is the reason.

\subsection{Topology as a measured object}
\label{sec:rel:measure}

A parallel line treats the interaction pattern as something to be observed
rather than designed. Oliveira, Bastos-Filho and Menezes reconstruct
\emph{influence graphs} from the information exchanges that actually occur
during a run and characterise swarms through network-science statistics,
including Laplacian spectra and component structure, using them to diagnose
stagnation and to compare
optimizers~\cite{oliveira2015network,oliveira2020uncovering}. Local optima
networks~\cite{ochoa2008lon} apply the same instinct to the landscape,
representing basins as nodes and search transitions as edges; as their authors
stress, the objective is characterisation, and the graph is built after the
fact rather than used to drive search.

The distinction matters for our purposes. Measurement studies establish that
graph statistics \emph{correlate} with search behaviour. They do not establish
that feeding such a statistic back into the update is beneficial, which is a
causal claim of a different order and the one we test.

\subsection{Theory of consensus-type updates}
\label{sec:rel:theory}

Consensus-based optimization (CBO) provides the most developed theory for
updates of the form \eqref{eq:master}. Particles drift toward a weighted
consensus point under multiplicative noise; convergence to global minimisers
has been established in the mean-field limit under conditions considerably
weaker than convexity~\cite{pinnau2017consensus,carrillo2018analytical}, and
subsequently for the finite-particle
scheme~\cite{fornasier2024cbo}. Polarized CBO localises the consensus through
a kernel, so that each particle is drawn toward a weighted mean emphasising
nearby particles~\cite{bungert2024polarized}; algebraically this is a
neighbourhood aggregate with a state-dependent, rather than graph-dependent,
weighting.

This body of work characterises what such operators \emph{do}: it identifies
the drift, quantifies the contraction of the empirical measure, and relates
the consensus rate to the operator's spectrum. It is, however, silent on the
question we pose. CBO theory takes the weighting as given and asks whether the
resulting dynamics converge; it does not ask whether one weighting is more
informative than another with the same distribution. Proposition~\ref{prop:mean}
is precisely a statement of that second kind, and to our knowledge has no
counterpart in this literature.

\subsection{Adaptive components and credit assignment}
\label{sec:rel:adaptive}

Choosing algorithm components online is a mature subfield. Adaptive operator
selection formulates the choice as a multi-armed bandit, assigning credit from
observed improvement and balancing exploration against exploitation with
rules ranging from probability matching and adaptive pursuit to
UCB variants~\cite{fialho2010analyzing}. Parameter control more broadly, and
hyper-heuristics, rest on the same premise: that feedback from the run
identifies the better alternative.

Outside metaheuristics, choosing a graph to optimise a spectral property is
routine---Ghosh and Boyd grow graphs to maximise algebraic
connectivity~\cite{ghosh2006growing}, and connectivity maintenance through the
Fiedler value is standard in multi-robot coordination. These methods optimise
a stated objective and their success is verifiable against it, which is not
the situation for an adaptive component inside a stochastic optimizer, where
the claim that selection is informed is seldom separated from the effect of
switching at all. Our second control (Section~\ref{sec:rand}) targets exactly
this gap, and we apply it to an adaptive scheme of our own construction as
well as to published ones.

\subsection{Methodological critique and ablation practice}
\label{sec:rel:method}

S\"orensen's critique of metaphor-driven algorithm design~\cite{sorensen2015metaheuristics}
prompted sustained attention to what novelty in this field consists of, and
subsequent programmatic work has pressed for component-level analysis in place
of monolithic algorithm proposals~\cite{swan2022metaheuristics}. Component-wise
ablation is now common in strong venues, and benchmark protocols such as
CEC2017~\cite{awad2016cec2017} have standardised the outcome measurements,
alongside guidance on statistical practice~\cite{derrac2011practical,garcia2009study},
sample sizing~\cite{campelo2019sample} and benchmarking more
broadly~\cite{bartz2020benchmarking}.

Ablation nonetheless answers a question adjacent to the one authors typically
ask. Deleting a component yields a different algorithm, with a different
operator spectrum and often a different runtime; the comparison therefore
confounds (P1) with (P2). Controls of the second kind---mechanism intact at
full strength, only the information it consumes destroyed---are standard
practice in fields that must separate an intervention from its delivery, and
permutation-based null models are long established for testing whether an
observed structure carries signal, from degree-preserving rewiring in network
biology~\cite{maslov2002specificity} to null models in community
ecology~\cite{gotelli2000null}. Comparable reality checks have recently
reshaped reported progress in neighbouring fields~\cite{musgrave2020metric,lipton2019troubling}. We are not
aware of their systematic use to audit structural attribution in
metaheuristics, and it is this transfer, together with the invariance
properties that make the transfer sound, that we contribute.
